The proof below uses the standard Connes–Consani correspondence between perfect difference sets and hyperfield extensions of the Krasner hyperfield, together with their description of 3-fold hyper-sums in the associated projective plane. ([math.jhu.edu](https://math.jhu.edu/~kc/papers/JNT2.pdf))

\documentclass[11pt]{article}
\usepackage{amsmath,amsthm,amssymb,mathtools}

\newtheorem{theorem}{Theorem}
\newtheorem{lemma}[theorem]{Lemma}
\newtheorem{remark}[theorem]{Remark}

\newcommand{\hz}{\mathbb{0}}

\begin{document}

\title{A Counterexample to a Cubic Identity over Hyperfields}
\author{}
\date{}
\maketitle

\begin{abstract}
We show that the implication
\[
\hz \in x\boxplus y\boxplus z
\quad\Longrightarrow\quad
\hz \in x^3\boxplus y^3\boxplus z^3\boxplus (-xyz)\boxplus(-xyz)\boxplus(-xyz)
\]
is \emph{false} for general hyperfields.
\end{abstract}

\section*{Statement}

We answer the following question.

\medskip

\noindent
\emph{If $\mathbb H$ is a hyperfield and $\hz \in x\boxplus y\boxplus z$, must one have}
\[
\hz \in x^3\boxplus y^3\boxplus z^3\boxplus (-xyz)\boxplus(-xyz)\boxplus(-xyz)\, ?
\]

\medskip

\noindent
The answer is \emph{no}.

\section*{A combinatorial lemma}

Recall that a subset $S\subseteq \mathbb Z$ is a \emph{Sidon set} if every nonzero integer has at most one representation as $s-s'$ with $s,s'\in S$. A subset $D\subseteq \mathbb Z$ is a \emph{perfect difference set} if every nonzero integer has exactly one such representation.

\begin{lemma}\label{lem:extend}
Every finite Sidon set $S\subseteq \mathbb Z$ is contained in an infinite perfect difference set.
\end{lemma}

\begin{proof}
We give the standard greedy construction.

Start with $D_0:=S$. Inductively, suppose $D_{n-1}$ is a finite Sidon set. If $n$ is already represented as a difference of two elements of $D_{n-1}$, set $D_n:=D_{n-1}$.

Assume now that $n$ is \emph{not} represented in $D_{n-1}$. We will add two new integers $N$ and $N+n$. Consider the finite family of affine functions
\[
\mathcal F_N:=
\{N-d,\ N+n-d,\ d-N,\ d-(N+n): d\in D_{n-1}\}\cup\{\pm n\}.
\]
If two distinct elements of $\mathcal F_N$ coincide, or if one of them coincides with a nonzero old difference from $D_{n-1}$, then $N$ satisfies a linear equation with integer coefficients; hence this gives only finitely many forbidden values of $N$. Also, requiring $N,N+n\notin D_{n-1}$ excludes only finitely many values.

Choose $N$ outside this finite forbidden set. Then all new nonzero differences created by adjoining $N$ and $N+n$ are pairwise distinct, are different from the old nonzero differences, and $(N+n)-N=n$ is the unique new representation of $n$. Hence
\[
D_n:=D_{n-1}\cup\{N,N+n\}
\]
is again a finite Sidon set.

Proceeding for $n=1,2,3,\dots$, we obtain an increasing sequence
\[
D_0\subseteq D_1\subseteq D_2\subseteq \cdots
\]
of finite Sidon sets, where after stage $n$ the integer $n$ is represented as a difference of two elements. Set
\[
D:=\bigcup_{n\ge 0} D_n.
\]
Because each $D_n$ is Sidon, every nonzero integer has \emph{at most} one representation as a difference of two elements of $D$; because at stage $n$ we ensured a representation of $n$, every positive integer has \emph{at least} one representation. Therefore every nonzero integer has exactly one representation, so $D$ is a perfect difference set containing $S$.
\end{proof}

\section*{Hyperfields from perfect difference sets}

We use the following standard construction (see Connes--Consani, \emph{The hyperring of ad\`ele classes}, Proposition~2.4, Proposition~3.1, Proposition~3.5, and Example~4.7).

\begin{lemma}\label{lem:CC}
Let $D\subseteq \mathbb Z$ be a perfect difference set. Then there exists a hyperfield
\[
\mathbb H_D=\{\hz\}\cup\{\tau^n:n\in\mathbb Z\}
\]
with the following properties.

\begin{enumerate}
    \item The multiplicative group of nonzero elements is infinite cyclic, generated by $\tau$:
    \[
    \tau^m\cdot \tau^n=\tau^{m+n}.
    \]
    \item Every nonzero element is self-opposite:
    \[
    -\tau^n=\tau^n,\qquad
    \tau^n\boxplus\tau^n=\{\hz,\tau^n\}.
    \]
    In particular,
    \[
    \tau^n\boxplus\tau^n\boxplus\tau^n=\{\hz,\tau^n\}.
    \]
    \item The lines in the associated projective plane are the translates $D+c$ of $D$.
    If $\tau^m\neq \tau^n$ lie on the unique line with exponent set $D+c$, then
    \[
    \tau^m\boxplus \tau^n=\{\tau^k:k\in (D+c)\setminus\{m,n\}\}.
    \]
    \item If $p,q,r$ are three distinct noncollinear nonzero points, then
    \[
    p\boxplus q\boxplus r
    \]
    is exactly the complement of the three sides of the triangle $pqr$.
\end{enumerate}
\end{lemma}

\begin{proof}
This is precisely the cited difference-set/projective-plane/hyperfield correspondence.
\end{proof}

\section*{The counterexample}

\begin{theorem}
The implication in the question is false in general.
\end{theorem}

\begin{proof}
Consider the finite set
\[
S:=\{0,2,7,20,26,29,200,212\}\subseteq \mathbb Z.
\]
Its positive differences are
\[
\begin{aligned}
&2,3,5,6,7,9,12,13,18,19,20,22,24,26,27,29,\\
&171,174,180,183,186,192,193,198,200,205,210,212,
\end{aligned}
\]
which are pairwise distinct. Hence $S$ is a Sidon set.

By Lemma~\ref{lem:extend}, there exists an infinite perfect difference set $D\subseteq \mathbb Z$ with
\[
S\subseteq D.
\]
Let $\mathbb H_D$ be the hyperfield from Lemma~\ref{lem:CC}. Set
\[
x:=\tau^0=1,\qquad y:=\tau^2,\qquad z:=\tau^7.
\]

\medskip

\noindent\textbf{Step 1: the hypothesis holds.}
Since $0,2,7\in D$, the points $\tau^0,\tau^2,\tau^7$ lie on the line with exponent set $D$.
By Lemma~\ref{lem:CC}(3),
\[
x\boxplus y=\{\tau^k:k\in D\setminus\{0,2\}\},
\]
so in particular $z=\tau^7\in x\boxplus y$. Because every nonzero element is self-opposite, this implies
\[
\hz\in x\boxplus y\boxplus z.
\]

\medskip

\noindent\textbf{Step 2: compute the cubic data.}
We have
\[
x^3=\tau^0,\qquad y^3=\tau^6,\qquad z^3=\tau^{21},\qquad xyz=\tau^{0+2+7}=\tau^9.
\]
Since all nonzero elements are self-opposite in $\mathbb H_D$,
\[
-xyz=xyz=\tau^9.
\]

\medskip

\noindent\textbf{Step 3: the point $\tau^9$ lies on the side joining $\tau^0$ and $\tau^6$.}
Because $20,26\in D$, the translate $D-20$ contains $0$ and $6$. Since $D$ is a perfect difference set, the representation
\[
6=26-20
\]
is unique; hence $D-20$ is the unique line through the points $\tau^0$ and $\tau^6$.

Now $29\in D$, so $9=29-20\in D-20$. Therefore
\[
\tau^9
\]
lies on the side joining $\tau^0$ and $\tau^6$.

\medskip

\noindent\textbf{Step 4: the three cube-points are noncollinear.}
If $\tau^{21}$ also lay on the line with exponent set $D-20$, then $41=21+20$ would belong to $D$. But then
\[
41-29=12=212-200,
\]
giving two distinct representations of the nonzero difference $12$ by elements of $D$, contradicting that $D$ is a perfect difference set. Hence
\[
\tau^{21}\notin D-20.
\]
So the three points
\[
\tau^0,\ \tau^6,\ \tau^{21}
\]
are distinct and noncollinear.

Let
\[
A:=x^3\boxplus y^3\boxplus z^3
=\tau^0\boxplus \tau^6\boxplus \tau^{21}.
\]
By Lemma~\ref{lem:CC}(4), $A$ is the complement of the three sides of the triangle with vertices $\tau^0,\tau^6,\tau^{21}$. Since $\tau^9$ lies on one of these sides, we get
\[
\tau^9\notin A.
\]
Also $\hz\notin A$.

\medskip

\noindent\textbf{Step 5: the required cubic hyper-sum does \emph{not} contain $\hz$.}
By Lemma~\ref{lem:CC}(2),
\[
\tau^9\boxplus\tau^9\boxplus\tau^9=\{\hz,\tau^9\}.
\]
Therefore
\[
\begin{aligned}
&x^3\boxplus y^3\boxplus z^3\boxplus(-xyz)\boxplus(-xyz)\boxplus(-xyz)\\
&\qquad = A\boxplus \tau^9\boxplus\tau^9\boxplus\tau^9\\
&\qquad = A\boxplus\{\hz,\tau^9\}\\
&\qquad = A\cup (A\boxplus \tau^9).
\end{aligned}
\]
Now $\hz\notin A$, and since every nonzero element is self-opposite,
\[
\hz\in A\boxplus \tau^9
\quad\Longleftrightarrow\quad
\tau^9\in A,
\]
which is false. Hence
\[
\hz\notin x^3\boxplus y^3\boxplus z^3\boxplus(-xyz)\boxplus(-xyz)\boxplus(-xyz).
\]

Thus we have exhibited a hyperfield $\mathbb H_D$ and elements $x,y,z\in \mathbb H_D$ such that
\[
\hz\in x\boxplus y\boxplus z
\]
but
\[
\hz\notin x^3\boxplus y^3\boxplus z^3\boxplus(-xyz)\boxplus(-xyz)\boxplus(-xyz).
\]
So the answer to the question is negative.
\end{proof}

\begin{thebibliography}{9}

\bibitem{CC}
A.~Connes and C.~Consani,
\emph{The hyperring of ad\`ele classes},
J. Number Theory \textbf{131} (2011), 159--194.

\end{thebibliography}

\end{document}